\chapter{Fundamentos estadísticos, aprendizaje y optimización del gemelo}
\label{ch:learning_foundations}

La calibración electromagnética y la predicción de sus discrepancias constituyen problemas
inversos relacionados, pero distintos. En el primero se ajustan parámetros del simulador;
en el segundo se aprende una función que aproxima un descriptor de error a partir de
información disponible en una consulta nueva. Este capítulo establece los fundamentos
necesarios para interpretar las regresiones de S4.4, los análisis agrupados de S4.3, la caracterización circular de S4.6 y la
extensión posterior a representaciones profundas. La formulación probabilística sigue
la distinción entre modelo, inferencia y decisión de Bishop, mientras el tratamiento de
redes profundas y regularización se apoya en Goodfellow et al.\ \citep{Bishop2006PRML,Goodfellow2016DL}.

\section{Aprendizaje supervisado y objetivo de generalización}

Sea $\vect z_i$ un descriptor disponible en inferencia y $y_i$ una variable objetivo.
Para S4.4, $\vect z_i$ contiene posición y/o descriptores del trazado nominal, mientras
$y_i$ es un retardo efectivo obtenido durante la construcción de las etiquetas a partir
del canal medido. Una función $g_\theta$ se estima mediante
\begin{equation}
\widehat\theta=\arg\min_\theta
\frac{1}{N_{\rm ent}}\sum_{i\in\mathcal I_{\rm ent}}
\ell\bigl(y_i,g_\theta(\vect z_i)\bigr)+\lambda\Omega(\theta).
\label{eq:empirical_risk}
\end{equation}
El riesgo de interés es, en cambio,
\begin{equation}
\mathcal R_Q(g)=\mathbb E_{(\vect z,y)\sim Q}
\left[\ell\bigl(y,g(\vect z)\bigr)\right],
\end{equation}
donde $Q$ representa la distribución de despliegue. Una partición aleatoria estima
principalmente un riesgo bajo condiciones semejantes a las observadas. Una región
reservada modifica la distribución de consultas y exige estudiar extrapolación espacial.
La separación entre ambos riesgos es central para interpretar S4.4d--f.

Construir las etiquetas con mediciones es compatible con aprendizaje supervisado.
La restricción consiste en impedir que el descriptor, la normalización o la selección
de hiperparámetros utilicen las mediciones de prueba. Cuando el descriptor proviene
de un simulador calibrado, sus parámetros también deben entrenarse sin las regiones
reservadas; alternativamente, debe declararse un protocolo con calibración previa común
y una afirmación de generalización condicionada a ella.

\section{Referencias de regresión y complejidad del predictor}

\subsection{Constante, vecino próximo y regresión regularizada}

Una constante ajustada por error absoluto utiliza la mediana de entrenamiento; por error
cuadrático utiliza la media. Esta elección debe distinguirse del uso arbitrario de cero.
El vecino próximo predice
\begin{equation}
\widehat y(\vect p)=y_{i^\star},\qquad
i^\star=\arg\min_{i\in\mathcal I_{\rm ent}}\|\vect p-\vect p_i\|_2.
\end{equation}
Es una referencia necesaria para determinar cuánto aporta un predictor frente a la
proximidad espacial. No estima un mecanismo electromagnético: interpola una etiqueta
en la geometría observada.

Para descriptores centrados y escalados con estadísticos de entrenamiento, la regresión
ridge resuelve
\begin{equation}
(\widehat b,\widehat{\vect w})=
\arg\min_{b,\vect w}\frac{1}{N_{\rm ent}}
\|\vect y-b\vect 1-\mat Z\vect w\|_2^2+\lambda\|\vect w\|_2^2.
\label{eq:ridge_model}
\end{equation}
Con datos centrados, $\widehat{\vect w}=(\mat Z^T\mat Z+
N_{\rm ent}\lambda\mat I)^{-1}\mat Z^T\vect y$ para $\lambda>0$.
El intercepto no se penaliza. La regularización estabiliza direcciones correlacionadas;
no corrige por sí sola una relación no lineal ni un cambio de dominio
\citep{Bishop2006PRML,Boyd2004Convex}.

\subsection{Bosques aleatorios y redes neuronales}

Un bosque de regresión combina árboles entrenados con aleatorización de muestras y
variables \citep{Breiman2001RandomForests}:
\begin{equation}
g_{\rm bosque}(\vect z)=\frac{1}{B}\sum_{b=1}^{B}T_b(\vect z).
\end{equation}
En las tablas experimentales se conserva la abreviatura RF del registro de ejecución
para \emph{Random Forest}; fuera de esas tablas RF designa radiofrecuencia. El método
puede representar interacciones no lineales entre descriptores sin exigir un modelo
paramétrico de su dependencia. En regresión convencional sus predicciones promedian
valores de las hojas, por lo que no debe suponerse extrapolación funcional ilimitada
fuera del soporte de entrenamiento. La importancia de una variable tampoco identifica
su efecto causal, especialmente cuando varios descriptores están correlacionados.

Una red multicapa compone transformaciones afines y no linealidades. Para materiales,
la salida puede parametrizarse mediante
\begin{equation}
\epsilon_r(\vect x)=1+\operatorname{softplus}(a_\psi(\vect x)),\qquad
\sigma(\vect x)=\operatorname{softplus}(b_\psi(\vect x)),
\label{eq:material_constraints}
\end{equation}
como una restricción suficiente para una clase de dieléctricos pasivos no magnéticos
considerados en el modelo. La cota $\epsilon_r\geq1$ es una elección de esa clase,
no una propiedad universal de todo medio electromagnético. Estas expresiones describen
una parametrización admisible; no atribuyen retrospectivamente esa implementación a B2.
Para una representación con frecuencia variable deben comprobarse además dispersión,
validez del modelo material y consistencia entre bandas.

La codificación posicional amplía las escalas espaciales representables, pero puede
facilitar la memorización. La capacidad, regularización y distribución de las mediciones
deben evaluarse conjuntamente. El número de parámetros, por sí solo, no determina la
fidelidad física ni la generalización \citep{Goodfellow2016DL}.

\section{Optimización diferenciable, restricciones e identificabilidad}

Para un objetivo real sobre canales complejos puede trabajarse con
$\vect u=[\Re H;\Im H]$. La regla de la cadena del \cref{ch:wdt} produce gradientes
respecto de parámetros reales sin tratar una pérdida no holomorfa como función analítica
compleja. La optimización del trazado se efectúa sobre un conjunto de caminos válido
en una región de parámetros; cambios de visibilidad pueden introducir discontinuidades.

Ridge con penalización positiva constituye un problema convexo; calibrar materiales,
geometría y redes neuronales conjuntamente no lo es en general. Por ello deben
distinguirse convergencia numérica, unicidad del mínimo e identificabilidad física
\citep{Boyd2004Convex,KaipioSomersalo2005}. Una pérdida pequeña puede coexistir con
parámetros diferentes que producen canales casi indistinguibles. Las múltiples
inicializaciones y la evaluación reservada caracterizan estabilidad empírica, pero no
demuestran unicidad global.

La pérdida de ganancia y dispersión de retardos de la calibración de referencia puede
formularse con un error relativo simétrico para cantidades no negativas,
\begin{equation}
\ell_{\rm sym}(a,b)=\frac{2|a-b|}{a+b+\varepsilon_a},\qquad
\mathcal L_{\rm cal}=\lambda_P\ell_{\rm sym}(P,\widehat P)
+\lambda_\tau\ell_{\rm sym}(\tau_{\rm RMS},\widehat\tau_{\rm RMS}).
\label{eq:calibration_macro_loss}
\end{equation}
Cada estabilizador debe tener las unidades de su denominador. Esta familia de objetivos
explica por qué mejorar magnitudes macroscópicas no impone concordancia de fase;
la forma y pesos exactos de una ejecución deben proceder de su configuración
\citep{Hoydis2024DiffRT}. Añadir términos complejos o diferenciales modifica el problema
de estimación y requiere una nueva comparación, con acceso a datos equivalente.

En optimización multiobjetivo, GradNorm ajusta pesos según magnitudes de gradiente y
ritmos de aprendizaje; PCGrad modifica componentes de gradientes conflictivos
\citep{Chen2018GradNorm,Yu2020PCGrad}. Ninguno garantiza una solución global, una
frontera de Pareto completa ni corrección del modelo físico. Su aplicación se condiciona
a observar un conflicto relevante entre objetivos previamente definidos.

\section{Estadística circular de los errores de fase}
\label{sec:circular_foundations}

Las diferencias angulares se definen en el círculo y no sobre una recta. Para residuos
$r_k=\wrap_{(-\pi,\pi]}(\phi_k-\widehat\phi_k)$, se consideran
\begin{equation}
\overline z=\frac{1}{K}\sum_k e^{jr_k},\qquad
\overline R=|\overline z|,\qquad V_{\rm circ}=1-\overline R.
\end{equation}
La dirección media es $\angle\overline z$ cuando la resultante permite definirla de
forma estable. Una media aritmética de ángulos próximos a $-\pi$ y $\pi$ puede ser
engañosa; las cantidades circulares evitan esa discontinuidad
\citep{Mardia2000Directional}.

La estructura espectral puede describirse mediante
\begin{equation}
C_\phi(L)=\left|\frac{1}{|\mathcal K_L|}
\sum_{k\in\mathcal K_L}e^{j(r_{k+L}-r_k)}\right|,
\label{eq:circular_lag}
\end{equation}
donde $\mathcal K_L$ contiene pares válidos separados por $L$ índices de frecuencia.
Esta resultante de incrementos no es una correlación de Pearson centrada: una fase
residual independiente pero concentrada también puede producir un valor positivo.
Para contrastar independencia se compara con un control de permutación que conserve
la distribución marginal y las máscaras. La concentración elevada a retardos cortos
indica regularidad local; no demuestra que el residual sea determinista ni identifica
su causa física. El ajuste previo de retardo y CPO puede introducir dependencia y
debe aplicarse igualmente al control.

La distribución de von Mises,
\begin{equation}
p(r\mid\mu,\kappa)=\frac{e^{\kappa\cos(r-\mu)}}{2\pi I_0(\kappa)},
\end{equation}
permite representar concentración circular. En una predicción de fase, las salidas
seno y coseno, o una distribución circular completa, evitan una pérdida euclídea
inadecuada cerca del corte angular. La concentración no se interpreta como incertidumbre
calibrada sin comprobar su cobertura sobre observaciones reservadas.

\section{Descriptores de trayectorias y dependencia entre niveles}
\label{sec:multilevel_foundations}

Para energías no negativas $e_\ell$ de caminos activos se define
$p_\ell=e_\ell/\sum_r e_r$ si la suma es positiva. La concentración de energía se
caracteriza con $D_1=\max p_\ell$, fracciones de los caminos dominantes,
$N_{\rm eff}=1/\sum p_\ell^2$ y entropía normalizada
$\mathcal H_p=-\sum p_\ell\log p_\ell/\log L$ para $L>1$.
Se usa $0\log0=0$; para un único camino se adopta entropía cero, y sin energía válida
el descriptor se marca como no definido. Estas probabilidades describen energías
individuales y no eliminan los términos de interferencia coherente del canal total.

Sean $x_{i,a}$ y $y_{i,a}$ descriptores de una posición $i$ y enlace $a$.
La asociación entre agregados $\overline x_i,\overline y_i$, la asociación de variables
centradas $x_{i,a}-\overline x_i$ y $y_{i,a}-\overline y_i$, y la distribución de
correlaciones calculadas por posición son estimandos distintos. Sus signos pueden
divergir sin contradicción. En S4.3 deben conservarse por separado las reducciones
utilizadas en cada tabla; agruparlas bajo un único coeficiente ocultaría la jerarquía.

El remuestreo agrupado selecciona posiciones con reemplazo y conserva conjuntamente
sus enlaces, volviendo a calcular el estadístico en cada réplica. Este procedimiento
respeta la dependencia interna del grupo \citep{Efron1993Bootstrap}; si posiciones
vecinas siguen correlacionadas, no elimina la dependencia entre grupos. La extensión
espacial requiere bloques o sesiones de remuestreo y un análisis de sensibilidad a
su tamaño. Un intervalo estrecho no transforma una asociación observacional en
identificación causal.

\section{Métricas, agregación y validación espacial}

Para errores escalares $e_i=\widehat y_i-y_i$ se distinguen
\begin{equation}
\mathrm{MAE}=N^{-1}\sum_i|e_i|,\quad
\mathrm{MedAE}=\operatorname{mediana}_i|e_i|,\quad
\mathrm{RMSE}=\sqrt{N^{-1}\sum_i e_i^2}.
\end{equation}
El coeficiente $R^2=1-\sum_i e_i^2/\sum_i(y_i-\overline y)^2$ puede ser negativo
y no está definido para una respuesta constante. Pearson mide asociación lineal y
Spearman asociación por rangos; ninguno garantiza acuerdo sin sesgo. La media de
medianas entre particiones difiere, en general, de la mediana de las predicciones
agrupadas. Esta diferencia explica reducciones distintas entre S4.4f y S4.5b.

En validación espacial, todos los enlaces de una posición pertenecen a la misma
partición. Los conjuntos de ajuste y prueba se separan mediante una región de exclusión
de anchura declarada. La selección de hiperparámetros utiliza sólo entrenamiento y
validación, preferiblemente con una separación interior compatible con la evaluación
exterior. Los cinco bloques de S4.4f caracterizan regiones de una escena; no constituyen
cinco edificios independientes ni demuestran transferencia entre dispositivos.

La predicción de un retardo auxiliar debe evaluarse también en la respuesta compleja
que modifica. Esa evaluación debe registrar qué parte de la corrección se predice
y cuál se estima con una medición disponible. Esta separación fundamenta el alcance
condicionado de S4.5 y prepara la posterior validación del operador de observación.


\section{Predictibilidad circular y alcance de los controles negativos}
\label{sec:circular_predictability}

Para una fase $\psi\in\mathbb S^1$ y variables disponibles $\vect z$, la pérdida
circular $\ell(\psi,g)=1-\cos(\psi-g)$ tiene riesgo condicional
\begin{equation}
\mathbb E[\ell(\psi,g)\mid\vect z]
=1-\Re\{e^{-jg}m(\vect z)\},\qquad
m(\vect z)=\mathbb E[e^{j\psi}\mid\vect z].
\end{equation}
Si $m(\vect z)\ne0$, el predictor óptimo para esta pérdida es
$g^\star(\vect z)=\angle m(\vect z)$ y el riesgo mínimo es $1-|m(\vect z)|$.
Si $m=0$, ninguna dirección puntual mejora ese riesgo condicional. Esta propiedad
no identifica automáticamente el predictor óptimo bajo distancia angular absoluta,
cuya solución corresponde a una mediana circular. Tampoco la resultante marginal
$|\mathbb E e^{j\psi}|$ determina la resultante condicional: fases predecibles desde
$\vect z$ pueden cancelarse al agregarlas \citep{Mardia2000Directional,Bishop2006PRML}.

Bajo el control específico $\Delta\psi\sim\mathcal U(-\pi,\pi]$, se cumple
\begin{equation}
\mathbb E|\Delta\psi|=\operatorname{Med}|\Delta\psi|=\frac{\pi}{2},\qquad
\mathbb E\cos\Delta\psi=0,\qquad
P(|\Delta\psi|\leq\alpha)=\frac{\alpha}{\pi},\quad 0\leq\alpha\leq\pi.
\label{eq:uniform_phase_control}
\end{equation}
Así, 90 grados, $1/6$ para 30 grados y $1/4$ para 45 grados son referencias
analíticas útiles para S4.6. Su proximidad empírica no prueba que el error sea
uniforme o independiente. El fracaso del vecino próximo descarta una ventaja
práctica clara de esa regla en el soporte evaluado; no establece que
$m(\vect z)=0$ para cualquier conjunto de descriptores.

La concentración finita también requiere cautela. Para $A$ fases independientes
uniformes, $\mathbb E\overline R^2=1/A$, aunque la resultante poblacional sea
nula. Por tanto, una dirección común obtenida de $A=32$ antenas puede aparecer
por ajuste a fluctuaciones. La comparación pertinente repite el estimador de
$\beta$, el centrado por arreglo y las máscaras sobre controles que conserven
las marginales, o estima la dirección con unas antenas y evalúa las restantes.
Las antenas reales no se suponen independientes por esta identidad de referencia.

\section{Estructura temporal con muestreo irregular y trazabilidad}
\label{sec:temporal_circular_protocol}

Para marcas temporales físicas $t_i$, un intervalo $B$ define los pares válidos
$\mathcal I_B=\{(i,j):i<j,\ |t_j-t_i|\in B\}$ dentro de un mismo arreglo.
Se separan tres estadísticos:
\begin{align}
D_\psi(B)&=\operatorname{Med}_{(i,j)\in\mathcal I_B}
\left|\wrap(\psi_j-\psi_i)\right|,\\
C_\psi(B)&=\frac{1}{|\mathcal I_B|}\sum_{(i,j)\in\mathcal I_B}
\cos(\psi_j-\psi_i),\\
R_\psi(B)&=\left|\frac{1}{|\mathcal I_B|}\sum_{(i,j)\in\mathcal I_B}
e^{j(\psi_j-\psi_i)}\right|.
\label{eq:temporal_phase_structure}
\end{align}
Una rotación determinista no nula puede producir $C_\psi$ pequeño y $R_\psi$
elevado; ambos indicadores responden a preguntas diferentes. Los intervalos
sin pares válidos no se rellenan con cero. Cada curva debe especificar número
de pares, separación efectiva, regla de selección y posibles restricciones
espaciales; usar todos los pares o una muestra de ellos define ponderaciones
diferentes. En S4.6d.2 se conserva la tabla reportada, sin atribuirle conteos
por intervalo que el registro suministrado no especifica.

Los pares que comparten una captura son dependientes. Los intervalos y las
comparaciones entre vecinos temporales, espaciales y aleatorios requieren
remuestreo por capturas, bloques o sesiones y réplicas del control aleatorio.
Se debe informar incertidumbre de la diferencia pareada; proximidad de medianas
sin un margen prefijado no constituye equivalencia estadística. La elección de
bloques temporales debe considerar el desplazamiento espacial del transmisor:
una campaña móvil no identifica por sí sola un proceso de deriva a posición fija.

Recuperar una marca temporal mediante igualdad de coordenadas es válido para
el subconjunto auditado sólo si la clave es única y la transformación de
coordenadas está documentada. En campañas con posiciones repetidas se necesita
un identificador original de captura. Deben preservarse archivo, índice de
origen, marca temporal, partición y versión de calibración; reproducir únicamente
una semilla de permutación no reconstruye necesariamente la secuencia histórica.

\section{Aprendizaje condicionado por repetibilidad y sensibilidad física}
\label{sec:repeatability_sensitivity_learning}

La función de aprendizaje se elegirá después de caracterizar la adquisición. Para una
representación $\mathcal V$, se compararán pares con estado físico fijo y distinta
realización instrumental, y pares con una perturbación física conocida:
\begin{align}
D_{\rm rep}&=d\bigl(\mathcal V(Y(\mathbf q,\eta_1)),
\mathcal V(Y(\mathbf q,\eta_2))\bigr),\nonumber\\
D_{\rm fis}(\delta\mathbf q)&=d\bigl(\mathcal V(Y(\mathbf q+\delta\mathbf q,\eta)),
\mathcal V(Y(\mathbf q,\eta))\bigr),\label{eq:repeat_physical_distances}\\
Q(\delta\mathbf q)&=\frac{\operatorname{med}D_{\rm fis}(\delta\mathbf q)}
{\operatorname{med}D_{\rm rep}+\epsilon_d}.
\label{eq:physical_repeatability_ratio}
\end{align}
Las condiciones instrumentales idénticas de la segunda expresión son un ideal
experimental; su discrepancia residual se medirá con referencias y repeticiones.
$Q$ es descriptivo y depende de la escala y de $\epsilon_d$. Un codificador constante
colapsa ambas distancias; por ello no se optimizará únicamente este cociente.
Se exigirán respuesta física no degenerada, incertidumbre y capacidad de discriminar
perturbaciones frente a pares nulos. Se informarán curvas de detección a falsa alarma
fijada con validación, errores de desplazamiento y resultados por orientación.

Las comparaciones iniciales incluirán magnitud, canal complejo, productos relativos,
autocodificador complejo, reconstrucción enmascarada y aprendizaje contrastivo temporal.
La rama latente física añadirá supervisión de desplazamiento, respuesta diferencial
y consistencia entre modalidades cuando estas referencias existan. Las pérdidas de
reconstrucción y de tarea se evaluarán por separado: reducir una no implica reducir
la otra. Las regularizaciones de varianza y las pruebas de recuperación del estado
permitirán detectar colapso, sin atribuir significado causal a coordenadas latentes
por inspección visual.

WiFi-JEPA proporciona una comparación específica de aprendizaje autosupervisado de
CSI para postura 3D \citep{Kim2026WiFiJEPA}; su objetivo publicado no demuestra
preservación de fase respiratoria. Se incorporará una vez validadas las entradas,
con presupuesto de datos y cómputo comparable y sin incluir sesiones de prueba en el
preentrenamiento. GradNorm o PCGrad se utilizarán sólo si el diagnóstico de gradientes
muestra desequilibrio o conflicto relevante; su incorporación no sustituirá una
hipótesis física. Las restricciones de Maxwell se limitarán inicialmente a casos
canónicos verificables, antes de considerar una PINN de habitación completa.
